\question
Кот Барсик учит хозяйку своему языку. В нём два звука: «мр» и «мяу». Каждому слову кошачьего языка Барсик сопоставляет код -- последовательность звуков; во фразе коды произносятся подряд, без пауз. Барсик хочет, чтобы любую фразу можно было расшифровать однозначно. Длины кодов уже придуманных слов: $2$, $3$, $3$, $4$ (сами коды ещё не выбраны). Какое наибольшее число новых слов с кодами длины не больше $5$ может добавить Барсик? Ответ обоснуйте.

---------------

Автор -- Дарья Гаркуша